\documentclass[11pt]{article}

% =====================================================
% Fonts and encoding
% =====================================================
\usepackage[T1]{fontenc}
\usepackage[utf8]{inputenc}
\usepackage{lmodern}

% =====================================================
% Page layout
% =====================================================
\usepackage{geometry}
\geometry{margin=1in}

% =====================================================
% Micro-typography
% =====================================================
\usepackage{microtype}
\UseMicrotypeSet[protrusion]{basicmath}
\microtypesetup{expansion=false}
\setlength{\emergencystretch}{3em}
\hyphenpenalty=10000
\exhyphenpenalty=10000

% =====================================================
% Graphics and tables
% =====================================================
\usepackage{graphicx}
\usepackage{booktabs}
\usepackage{tabularx}
\usepackage{xcolor}

% =====================================================
% Mathematics
% =====================================================
\usepackage{amsmath, amssymb, amsthm}
\usepackage{mathtools}

% =====================================================
% Algorithms
% =====================================================
\usepackage{algorithm}
\usepackage{algorithmic}

% =====================================================
% References
% =====================================================
\usepackage[numbers,sort&compress]{natbib}

% =====================================================
% Hyperlinks
% =====================================================
\usepackage[hidelinks]{hyperref}

% =====================================================
% Paragraph formatting
% =====================================================
\setlength{\parskip}{0.3em}
\setlength{\parindent}{0pt}

% =====================================================
% Custom commands
% =====================================================
\usepackage{xspace}
\newcommand{\HCO}{HCO\xspace}

% =====================================================
% =====================================================
% Theorem environments
% =====================================================
\newtheorem{theorem}{Theorem}
\newtheorem{lemma}{Lemma}
\newtheorem{corollary}{Corollary}
\newtheorem{definition}{Definition}
\newtheorem{assumption}{Assumption}
\newtheorem{proposition}{Proposition}
\theoremstyle{remark}         
\newtheorem{remark}{Remark}    
\theoremstyle{plain}          
% =====================================================
% Title
% =====================================================
\title{Human Challenge Oracle: A Cryptographic Primitive\\
	for Continuous, Identity-Bound Sybil Resistance}

\author{Homayoun Maleki\\
	DeustoTech, University of Deusto, Bilbao, Spain}

\date{}

\begin{document}
	
	\maketitle
	
	\begin{abstract}
		Open systems face persistent Sybil attacks in which adversaries sustain large
		numbers of fake identities at sublinear cost. Existing defenses---CAPTCHAs,
		behavioral biometrics, and Proof-of-Personhood mechanisms---either verify
		identity only at creation time or impose costs that can be reused and amortized
		across identities. Neither paradigm enforces a cost that grows proportionally
		with both the number of sustained identities and the duration of participation.
		
		We introduce the \emph{Human Challenge Oracle}~(\HCO), a cryptographic
		primitive for continuous, identity-bound, rate-limited verification. \HCO\
		issues fresh challenges that are cryptographically bound to individual
		identities and time windows; valid responses must arrive within a strict
		deadline. The security argument rests on four structural
		properties---identity binding~(H1), freshness~(H2), real-time
		constraint~(H3), and per-channel throughput bound~(H4)---and is
		\emph{agnostic to solver identity}: whether a participation channel is
		operated by a human or an automated agent, the cost of sustaining $s$
		identities over $T$ windows is $\Omega(sT)$. No assumption about human
		cognitive superiority over automated systems is required.
		
		We present a concrete hash-based construction satisfying~(H1)--(H3)
		under standard cryptographic assumptions, with a formally specified
		verification protocol. We characterize four abstract challenge families
		admissible under~(H1)--(H4) and describe practical browser-based
		instantiations. A preliminary empirical study validates the feasibility
		of the throughput bound across representative families under strict
		response deadlines against current vision-language models. \HCO\ is
		designed as a composable primitive; its formal use as a consensus
		resource is developed in companion
		papers~\cite{maleki2026pocmt,maleki2026geometry}.
	\end{abstract}
	
	
	%% =============================================================
	\section{Introduction}
	\label{sec:intro}
	%% =============================================================
	
	Open and permissionless systems assume, implicitly or explicitly, that
	each account represents an independent participant. This assumption is
	violated by Sybil attacks~\cite{douceur2002sybil}, in which a single
	adversary creates and operates many fake identities at negligible cost,
	enabling spam, vote manipulation, coordinated disinformation, and
	governance abuse. Preventing Sybil attacks at scale remains one of
	the central unsolved problems in distributed systems security.
	
	\noindent\textbf{The failure mode of existing defenses.}
	The dominant approaches bind participation to a scarce resource.
	CAPTCHAs~\cite{vonahn2003captcha} exploit a perceptual gap between
	humans and automated agents at account creation time, but they are
	fundamentally one-shot: once an account is created, sustained Sybil
	participation incurs no proportional recurring cost. As modern
	vision-language models (VLMs) improve, this gap steadily
	erodes~\cite{sivakorn2016iam,ye2018yet,goodfellow2014multidigit}.
	CAPTCHA-solving services further commoditize the
	bypass~\cite{motoyama2010recaptchas}. Text-based CAPTCHA mechanisms
	have been shown to be breakable at high rates even against hardened
	designs~\cite{bursztein2011text}.
	
	Proof-of-Personhood (PoP) systems~\cite{borge2017proofofpersonhood,ford2020pseudonym}
	enforce a one-person--one-identity mapping through ceremonies, biometric
	enrollment, or social vouching. Deployed systems include
	Worldcoin~\cite{worldcoin2023}, BrightID~\cite{brightid2022}, and
	Gitcoin Passport~\cite{gitcoin2023}. Like CAPTCHAs, PoP mechanisms
	raise the cost of identity \emph{creation} but do not constrain ongoing
	participation: once identities are admitted, adversaries may sustain
	many of them without proportional recurring cost.
	
	Social-graph defenses~\cite{yu2006sybilguard,yu2008sybillimit,danezis2009sybilinfer,wang2013clickstream}
	constrain Sybil proliferation through trust topology but provide no
	explicit per-identity per-window cost bound. Behavioral
	biometrics~\cite{frank2013touchalytics,serwadda2013kids} enable ongoing
	verification but are proprietary and not designed for open adversarial
	settings.
	
	\noindent\textbf{The structural gap.}
	All of these defenses fail to enforce a cost that scales
	\emph{jointly} with the number of sustained identities $s$ and the
	participation horizon $T$. An adversary sustaining $s$ identities over
	$T$ windows should face a cost of $\Omega(sT)$---not a one-time fee at
	creation, and not a resource that can be reused across identities once
	acquired. The impossibility of achieving this under parallelizable
	resources (computation, capital) is established in the companion
	axiomatic paper~\cite{maleki2026geometry}.
	
	\noindent\textbf{This paper: \HCO\ as a cryptographic primitive.}
	We introduce the \emph{Human Challenge Oracle}~(\HCO), a security
	primitive designed to enforce exactly this cost structure. \HCO\
	repeatedly issues fresh, cryptographically identity-bound challenges
	that must be answered within a strict real-time deadline. The security
	argument depends on four structural properties:
	
	\begin{itemize}
		\item \textbf{(H1) Identity binding.} A valid response for identity
		$v$ cannot be credited to any other identity $v' \neq v$.
		\item \textbf{(H2) Freshness.} Challenges are unpredictable before
		issuance; precomputation is structurally impossible.
		\item \textbf{(H3) Real-time constraint.} Responses arriving after
		deadline $\Delta_{\mathrm{resp}}$ are rejected.
		\item \textbf{(H4) Per-channel throughput bound.} Any single
		participation channel---human or automated---produces at most
		$\tau_h = O(1)$ valid responses per window.
	\end{itemize}
	
	Crucially, the linear cost guarantee does \emph{not} assume that
	humans outperform automated solvers on individual challenges.
	Whether a channel is operated by a human, an AI system, or any hybrid
	is orthogonal to the cost argument: what cannot be avoided is that each
	identity requires one independent channel-window, and channel-windows
	cannot be shared across identities or carried across time.
	
	This solver-agnostic framing aligns \HCO\ with its companion papers.
	The \HCO\ primitive is used as the participation resource in the
	PoCmt consensus protocol~\cite{maleki2026pocmt}, which proves
	backbone-style safety and liveness under partial synchrony. The
	axiomatic theory establishing that linear Sybil cost requires a
	non-parallelizable resource is developed in~\cite{maleki2026geometry}.
	
	\noindent\textbf{Contributions.}
	\begin{enumerate}
		\item We formally specify \HCO\ as a cryptographic oracle
		primitive with properties~(H1)--(H4)~(Section~\ref{sec:spec}).
		\item We present a hash-based construction satisfying~(H1)--(H3)
		under standard assumptions, with a verified algorithmic
		specification~(Section~\ref{sec:construction}).
		\item We prove that (H1)--(H4) enforce $C(s,T) = \Omega(sT)$,
		both per window and in steady state, solver-agnostically
		(Sections~\ref{sec:cost}--\ref{sec:multiwindow}).
		\item We characterize four admissible challenge families and
		describe browser-based instantiations~(Section~\ref{sec:challenges}).
		\item We present a preliminary empirical evaluation of the throughput
		bound~(H4) against current VLMs under strict deadlines
		(Section~\ref{sec:eval}).
		\item We analyze the construction against relay, outsourcing,
		replay, precomputation, deepfake, and side-channel
		attacks~(Section~\ref{sec:security}).
	\end{enumerate}
	
	\noindent\textbf{\HCO\ is not rate limiting.}
	Standard rate-limiting mechanisms bound request frequency per account
	or network endpoint. They share none of the structural properties that
	make \HCO\ a formal cryptographic primitive:
	\begin{itemize}
		\item \emph{Rate limits do not enforce identity binding~(H1).}
		An adversary creates additional accounts; each gets its own rate
		budget. The Sybil cost of creating $s$ accounts remains
		zero---the rate limit is per-account, not per-channel-window.
		\item \emph{Rate limits do not enforce freshness~(H2).}
		An adversary precomputes, stockpiles, or replays requests across
		time windows without penalty.
		\item \emph{Rate limits do not yield a composable cost theorem.}
		There is no equivalent of $C(s,T) = \Omega(sT)$ for sustained
		multi-identity participation over time: without~(H1) and~(H2),
		the cost of identity replication does not compound.
		\item \emph{Rate limits do not compose into higher-level protocols.}
		They produce no cryptographically verifiable, identity-bound
		participation record usable as a weighting resource in a consensus
		mechanism or governance system.
	\end{itemize}
	The contribution of \HCO\ is the \emph{unification} of
	cryptographic identity binding, freshness, window locality, and
	throughput boundedness into a single formal primitive that admits a
	provable, composable linear cost theorem. Engineering mechanisms
	such as session serialization or device attestation are concrete
	instantiations of Property~(H4)---they are how~(H4) is
	\emph{deployed}, not what \HCO\ \emph{is}.
	
	
	%% =============================================================
	\section{Related Work}
	\label{sec:related}
	%% =============================================================
	
	\noindent\textbf{Sybil attacks and graph-based defenses.}
	Douceur~\cite{douceur2002sybil} established that identity replication
	cannot be prevented without binding influence to a scarce resource.
	Graph-based approaches constrain Sybil regions using social network
	structure: SybilGuard~\cite{yu2006sybilguard} and
	SybilLimit~\cite{yu2008sybillimit} bound adversarial identity count
	through random-walk protocols; SybilInfer~\cite{danezis2009sybilinfer}
	applies probabilistic inference over trust graphs. Clickstream
	behavioral analysis detects Sybil accounts through usage
	patterns~\cite{wang2013clickstream}. These mechanisms constrain the
	\emph{number} of Sybil identities but provide no explicit per-identity
	per-window cost bound over time.
	
	\noindent\textbf{CAPTCHA and one-time verification.}
	The CAPTCHA framework~\cite{vonahn2003captcha} exploits human
	perceptual advantage at account creation. reCAPTCHA~\cite{recaptcha2019}
	and hCaptcha~\cite{hcaptcha2020} extend this with behavioral signals and
	ML-based risk scoring. These mechanisms are fundamentally one-time: the
	cost of sustained Sybil participation after account creation is zero.
	Deep learning now breaks text~\cite{bursztein2011text,ye2018yet} and
	semantic image CAPTCHAs~\cite{sivakorn2016iam} at high rates, including
	via multi-digit recognition from distorted
	imagery~\cite{goodfellow2014multidigit}. CAPTCHA-solving services
	automate the bypass at commodity prices~\cite{motoyama2010recaptchas}.
	
	\noindent\textbf{Proof-of-Personhood.}
	PoP systems~\cite{borge2017proofofpersonhood,ford2020pseudonym} enforce
	a one-human--one-identity mapping through physical ceremonies or social
	vouching. Worldcoin~\cite{worldcoin2023}, BrightID~\cite{brightid2022},
	and Gitcoin Passport~\cite{gitcoin2023} illustrate deployed approaches.
	PoP mechanisms raise identity creation cost but provide no continuous
	guarantee; once admitted, adversaries can sustain many identities
	without proportional recurring cost.
	
	\noindent\textbf{Human--AI performance under time constraints.}
	Humans exhibit perceptual generalization under noise and distribution
	shift that deep networks lack~\cite{geirhos2018generalisation,geirhos2019texture},
	and shortcut learning limits model robustness in novel
	conditions~\cite{geirhos2020shortcut}. Broader cognitive flexibility
	distinguishes human reasoning from current AI
	systems~\cite{lake2017building}. These findings support the empirical
	plausibility of~(H4) for human channels, though our security argument
	does not depend on this gap being permanent.
	
	\noindent\textbf{Automated solver capabilities.}
	State-of-the-art VLMs---GPT-4o~\cite{openai2023gpt4},
	Gemini~\cite{geminiteam2023}, Claude~\cite{anthropic2024claude}---achieve
	strong performance on unconstrained visual and reasoning tasks.
	CLIP~\cite{radford2021clip} demonstrates zero-shot visual recognition.
	However, these models incur inference latencies of several seconds,
	which conflicts with tight $\Delta_{\mathrm{resp}}$ windows. Our
	evaluation tests these models under strict deadline enforcement.
	
	\noindent\textbf{Behavioral biometrics and continuous authentication.}
	Touchalytics~\cite{frank2013touchalytics} demonstrates touchscreen
	behavior as a continuous authentication signal; adversarial impersonation
	exposes fundamental limits of behavioral
	biometrics~\cite{serwadda2013kids}. These mechanisms are passive and
	proprietary, not designed for open adversarial settings or
	permissionless participation.
	
	\noindent\textbf{Trusted hardware.}
	Intel SGX~\cite{costan2016sgx} and related environments bind computation
	to physical hardware, enforcing per-device throughput limits. \HCO\
	does not require trusted hardware; it relies on cryptographic binding
	and the structural properties of participation channels.
	
	\noindent\textbf{Resource-based consensus.}
	PoW~\cite{nakamoto2008bitcoin} and PoS~\cite{kiayias2017ouroboros} bind
	influence to parallelizable resources that can be subdivided and reused
	across identities at sublinear marginal cost~\cite{bonneau2015sok}.
	The formal impossibility of linear Sybil cost under parallelizable
	resources is established in~\cite{maleki2026geometry}; the consensus
	protocol that uses \HCO\ to escape this impossibility is
	in~\cite{maleki2026pocmt}.
	
	\begin{table*}[t]
		\centering
		\caption{Comparison of verification mechanisms on dimensions critical
			for sustained Sybil resistance.}
		\label{tab:hco-taxonomy}
		\begin{tabular}{lcccc}
			\toprule
			Mechanism & Persistent & Per-identity Rate-Limited
			& Solver-Agnostic & Cost Scaling \\
			\midrule
			CAPTCHAs~\cite{vonahn2003captcha}
			& No  & No  & No  & One-time \\
			reCAPTCHA / hCaptcha~\cite{recaptcha2019,hcaptcha2020}
			& No  & No  & No  & One-time \\
			Proof-of-Personhood~\cite{borge2017proofofpersonhood}
			& No  & No  & No  & One-time \\
			Behavioral biometrics~\cite{frank2013touchalytics}
			& Yes & Yes & Partial & Unclear \\
			Proof-of-Work~\cite{nakamoto2008bitcoin}
			& Yes & No  & Yes & $o(sT)$ \\
			Proof-of-Stake~\cite{kiayias2017ouroboros}
			& Yes & No  & Yes & $o(sT)$ \\
			\textbf{\HCO\ (this work)}
			& Yes & Yes & Yes & $\Omega(sT)$ \\
			\bottomrule
		\end{tabular}
	\end{table*}
	
	Table~\ref{tab:hco-taxonomy} positions \HCO\ relative to prior work.
	\HCO\ is the only mechanism that simultaneously enforces persistence,
	per-identity rate limiting, solver-agnosticism, and linear joint cost.
	
	
	%% =============================================================
	\section{HCO Specification}
	\label{sec:spec}
	%% =============================================================
	
	\subsection{Oracle Interface}
	
	\begin{definition}[Human Challenge Oracle]
		\label{def:hco}
		A \emph{Human Challenge Oracle} is a stateful oracle
		\[
		\HCO(v,\, d,\, j) \;\longrightarrow\; \chi_{v,d,j}
		\]
		that, on input a validator identity $v$, window index $d \in \mathbb{N}$,
		and challenge index $j \in \{1,\ldots,k(d)\}$, outputs a fresh challenge
		instance $\chi_{v,d,j}$. Each challenge is associated with a response
		deadline $\Delta_{\mathrm{resp}} \ll \Delta_h$, where $\Delta_h > 0$ is
		the window duration.
	\end{definition}
	
	A response $\rho$ to $\chi_{v,d,j}$ is \emph{valid} iff:
	(i)~$\rho$ is received within $\Delta_{\mathrm{resp}}$ of issuance;
	(ii)~$\rho$ is a correct solution to $\chi_{v,d,j}$;
	(iii)~$\rho$ is cryptographically bound to the triple $(v,d,j)$.
	Verification is deterministic and publicly executable without long-term
	storage of participant data.
	
	\subsection{Security Properties}
	
	\HCO\ must satisfy the following four structural properties.
	
	\begin{itemize}
		\item \textbf{(H1) Identity binding.} Each challenge $\chi_{v,d,j}$
		is cryptographically bound to the triple $(v,d,j)$. A valid response
		to $\chi_{v,d,j}$ is invalid for any other triple
		$(v',d',j') \neq (v,d,j)$. Responses cannot be shared across
		identities or reused across windows.
		
		\item \textbf{(H2) Freshness.} No computationally bounded adversary
		can produce a valid response to $\chi_{v,d,j}$ with non-negligible
		probability before $\chi_{v,d,j}$ is issued. Precomputation and
		stockpiling are structurally impossible.
		
		\item \textbf{(H3) Real-time constraint.} Responses submitted after
		deadline $\Delta_{\mathrm{resp}}$ are rejected. Together with~(H2),
		this rules out any strategy separating challenge receipt from response
		generation across time.
		
		\item \textbf{(H4) Per-channel throughput bound.} There exists a
		constant $\tau_h = O(1)$ such that any single participation
		channel---human or automated---can produce valid responses for at most
		$\tau_h$ distinct challenge indices per window, independently of the
		number of identities associated with that channel.
	\end{itemize}
	
	\begin{proposition}[Security is solver-agnostic]
		\label{prop:agnostic}
		Under~(H1)--(H4), sustaining $s$ active identities over $T$ windows
		requires $\Omega(sT)$ independent channel-windows, regardless of
		whether channels are operated by humans, automated agents, or any
		combination thereof.
	\end{proposition}
	
	The formal proof appears in Sections~\ref{sec:cost}--\ref{sec:multiwindow}.
	
	
	\begin{remark}
		Properties~(H1)--(H4) jointly constitute a strictly stronger primitive
		than per-account rate limiting. Rate limiting bounds request frequency
		but provides no cryptographic identity binding, no freshness guarantee,
		and no composable cost theorem. \HCO\ provides all four, enabling the
		$C(s,T) = \Omega(sT)$ result of
		Sections~\ref{sec:cost}--\ref{sec:multiwindow} and composability with
		higher-level protocols such as the consensus mechanism
		of~\cite{maleki2026pocmt}.
	\end{remark}
	
	\subsection{On Property~(H4) and Solver Identity}
	
	Property~(H4) asserts that a single channel is throughput-bounded per
	window under $\Delta_{\mathrm{resp}}$. This is \emph{not} equivalent to
	assuming that humans outperform machines on individual challenges. A single
	automated agent may solve any given challenge instantly; what~(H4) asserts
	is that a single channel cannot handle $\omega(1)$ distinct challenge
	indices per window within $\Delta_{\mathrm{resp}}$, due to the irreducible
	latency of each challenge-response round trip (network delivery, processing,
	submission, and verification). For human channels, cognitive processing
	time widens this bound further. Empirical validation of~(H4) for specific
	challenge families is given in Section~\ref{sec:eval}.
	
	\noindent\textbf{Two independent bases for~(H4).}
	Property~(H4) can be grounded in two structurally distinct mechanisms,
	and either one is sufficient for the security argument.
	
	\emph{Cognitive serialization.}
	Human cognition imposes an irreducible serial bottleneck: a single person
	cannot attend to, process, and respond to multiple independent challenges
	simultaneously within a short deadline. This bound derives from cognitive
	and motor throughput limits, not from any assumption of human superiority
	over automated solvers.
	
	\emph{Channel serialization.}
	Engineering constraints can enforce~(H4) independently of who or what
	operates the channel. Concrete mechanisms include:
	\begin{itemize}
		\item \textbf{Per-session serialization.} The system issues at most
		one active challenge per session token at any moment; the browser
		session or WebSocket connection is the channel boundary. A channel
		serving multiple identities in parallel must hold multiple concurrent
		sessions, each incurring its own round-trip budget.
		\item \textbf{Device-bound attestation.} Challenges are bound to
		a specific hardware token, secure enclave, or TPM that enforces
		a per-device rate limit. Scaling to $s$ identities requires $s$
		independently attested devices.
		\item \textbf{Stateful interaction locality.} Challenges require
		continuous interaction with a dynamically mutating UI element
		throughout $\Delta_{\mathrm{resp}}$, making parallel execution across
		identities on a single device structurally difficult.
	\end{itemize}
	
	Under channel serialization, the bound on $\tau_h$ is a systems
	constraint rather than a cognitive claim: it follows from the
	serialization enforced by the channel architecture. The formal security
	argument of Sections~\ref{sec:cost}--\ref{sec:multiwindow} holds under
	either basis for~(H4); deployments may select the enforcement mechanism
	best suited to their trust model and accessibility requirements.
	
	\subsection{Separability and Composability}
	
	The security guarantees of \HCO\ depend only on~(H1)--(H4). Application-level
	details---user interfaces, challenge formats, delivery channels---do not affect
	these guarantees. This separability enables \HCO\ to be instantiated across
	diverse systems and composed with higher-level protocols. Its formal use as a
	primitive in the PoCmt consensus mechanism is developed
	in~\cite{maleki2026pocmt}.
	
	
	%% =============================================================
	\section{Cryptographic Construction}
	\label{sec:construction}
	%% =============================================================
	
	We present a concrete cryptographic instantiation of \HCO\ satisfying
	Properties~(H1)--(H3) under standard assumptions. Property~(H4)
	is analyzed separately as a structural and empirical claim in
	Sections~\ref{sec:spec} and~\ref{sec:eval}.
	
	\subsection{Setup}
	
	Let $\mathcal{H}: \{0,1\}^* \to \{0,1\}^\lambda$ be a collision-resistant
	hash function modeled as a random oracle~\cite{bellare1993random}, with
	security parameter $\lambda$. We instantiate $\mathcal{H}$ with
	SHA-3~\cite{nist2015sha3} in practice. Each identity $v$ holds a
	long-term key pair $(\mathsf{sk}_v, \mathsf{pk}_v)$ with $\mathsf{pk}_v$
	registered with the system.
	
	At the start of each window $d$, a fresh nonce
	$\eta_d \in \{0,1\}^\lambda$ is drawn from a public randomness
	beacon~\cite{syta2017drand} and publicly broadcast. The nonce is
	unpredictable before window $d$ opens and publicly verifiable thereafter.
	
	\subsection{Challenge Generation}
	
	For identity $v$, window $d$, and index $j \in \{1,\ldots,k(d)\}$:
	\begin{equation}
		\chi_{v,d,j}
		\;=\;
		\Bigl(
		\underbrace{\mathcal{H}(\mathsf{pk}_v \;\|\; d \;\|\; j \;\|\; \eta_d)}_{\text{binding tag}},
		\;\;
		\underbrace{\pi_{v,d,j}}_{\text{task payload}}
		\Bigr),
		\label{eq:challenge-gen}
	\end{equation}
	where $\pi_{v,d,j}$ is a task payload from an admissible challenge
	family (Section~\ref{sec:challenges}). The binding tag anchors the
	challenge to $(v,d,j)$ and the window-specific nonce.
	
	\subsection{Response and Verification}
	
	A participant submits:
	\[
	\rho \;=\; \bigl(\sigma_v,\; r_{v,d,j}\bigr),
	\]
	where $r_{v,d,j} \in \mathcal{R}$ is the task solution and
	$\sigma_v = \mathsf{Sign}(\mathsf{sk}_v,\;
	\mathsf{pk}_v \;\|\; d \;\|\; j \;\|\; r_{v,d,j})$.
	Algorithm~\ref{alg:verify} gives the complete verification procedure.
	
	\begin{algorithm}[t]
		\caption{$\HCO.\mathsf{Verify}(v,\, d,\, j,\, \rho,\, t_{\mathrm{recv}})$}
		\label{alg:verify}
		\begin{algorithmic}[1]
			\REQUIRE Identity $v$, window $d$, index $j$,
			response $\rho = (\sigma_v, r_{v,d,j})$,
			receipt time $t_{\mathrm{recv}}$
			\ENSURE \textbf{Accept} or \textbf{Reject}
			\IF{$t_{\mathrm{recv}} > t_{\mathrm{issue}} + \Delta_{\mathrm{resp}}$}
			\RETURN \textbf{Reject} \hfill\COMMENT{Deadline exceeded --- enforces (H3)}
			\ENDIF
			\IF{$\lnot\;\mathsf{VerifySig}\!\bigl(\mathsf{pk}_v,\;
				\mathsf{pk}_v \;\|\; d \;\|\; j \;\|\; r_{v,d,j},\; \sigma_v\bigr)$}
			\RETURN \textbf{Reject} \hfill\COMMENT{Invalid signature --- enforces (H1)}
			\ENDIF
			\STATE Recompute binding tag:
			$\tau \leftarrow \mathcal{H}(\mathsf{pk}_v \;\|\; d \;\|\; j \;\|\; \eta_d)$
			\IF{binding tag in $\chi_{v,d,j}$ $\neq$ $\tau$}
			\RETURN \textbf{Reject} \hfill\COMMENT{Wrong identity/window --- enforces (H1)}
			\ENDIF
			\IF{$\lnot\;\mathsf{CheckTask}(r_{v,d,j},\; \pi_{v,d,j})$}
			\RETURN \textbf{Reject} \hfill\COMMENT{Incorrect task solution}
			\ENDIF
			\RETURN \textbf{Accept}
		\end{algorithmic}
	\end{algorithm}
	
	\subsection{Security of the Construction}
	
	\noindent\textbf{Identity binding~(H1).}
	The binding tag $\mathcal{H}(\mathsf{pk}_v \;\|\; d \;\|\; j \;\|\; \eta_d)$
	ties the challenge to $(v,d,j)$. A valid response includes a signature
	under $\mathsf{sk}_v$; producing a valid response for $v' \neq v$ requires
	forging a signature under $\mathsf{sk}_v$ or finding a hash collision, both
	infeasible under standard assumptions.
	
	\noindent\textbf{Freshness~(H2).}
	The nonce $\eta_d$ is revealed only at the start of window $d$.
	Because $\eta_d$ is hidden prior to window $d$, the binding tag is
	computationally indistinguishable from uniform before $d$ opens. No
	adversary can compute $\chi_{v,d,j}$ before $\eta_d$ is known,
	ruling out precomputation in the random oracle model~\cite{bellare1993random}.
	
	\noindent\textbf{Real-time constraint~(H3).}
	Line~1 of Algorithm~\ref{alg:verify} checks the receipt timestamp
	deterministically and independently of the task solution. Late responses
	are rejected unconditionally.
	
	\noindent\textbf{Throughput bound~(H4).}
	The cryptographic construction does not by itself enforce~(H4); this
	property arises from the latency structure of the participation channel.
	Each valid response requires: (i)~receiving $\chi_{v,d,j}$ after its
	issuance; (ii)~computing $r_{v,d,j}$ and $\sigma_v$; (iii)~transmitting
	$\rho$ and having it received within $\Delta_{\mathrm{resp}}$. The minimum
	round-trip latency $\delta_{\mathrm{rtt}}$ limits any channel to at most
	$\tau_h = \lfloor \Delta_{\mathrm{resp}} / \delta_{\mathrm{rtt}} \rfloor$
	responses per window. For human channels, cognitive processing time further
	reduces this ceiling. Empirical characterization of $\tau_h$ is given
	in Section~\ref{sec:eval}.
	
	
	%% =============================================================
	\section{System and Adversary Model}
	\label{sec:model}
	%% =============================================================
	
	\subsection{System Model}
	
	The system consists of a set of identities $\mathcal{I}$ interacting with
	a platform relying on \HCO\ for continuous verification. Time is divided
	into windows $d \in \mathbb{N}$ of fixed duration $\Delta_h$.
	
	An identity $v$ is \emph{active} in window $d$ iff it submits at least
	one valid \HCO\ response during $d$. Verification is deterministic,
	public, and requires no long-term storage of participant data.
	
	\subsection{Adversary Model}
	
	We consider a computationally unbounded adversary $\mathcal{A}$ that may:
	\begin{itemize}
		\item create and coordinate an arbitrary number of identities $s$;
		\item operate up to $m$ independent participation channels (human or
		automated) within any window, with $m$ adaptively varying across windows;
		\item employ outsourcing, relay attacks, AI automation, or any combination;
		\item deviate arbitrarily from the protocol, subject to~(H1)--(H4).
	\end{itemize}
	
	\subsection{Adversarial Cost Function}
	
	\begin{definition}[Adversarial cost]
		$C(s, T)$ denotes the minimum number of independent participation
		channel-windows required to sustain $s$ active identities over $T$
		consecutive windows.
	\end{definition}
	
	
	%% =============================================================
	\section{Sybil Cost Analysis}
	\label{sec:cost}
	%% =============================================================
	
	\begin{lemma}[Per-window identity bound]
		\label{lem:identity-bound}
		In any window $d$, an adversary operating $m$ participation channels
		can produce valid \HCO\ responses for at most $m \cdot \tau_h$
		distinct identities.
	\end{lemma}
	
	\begin{proof}
		By~(H4), each of $m$ channels produces at most $\tau_h$ valid
		responses per window. By~(H1), a response for identity $v$ cannot
		be credited to $v' \neq v$; responses are non-sharable across
		identities. By~(H2) and~(H3), precomputed or late responses are
		rejected. The total valid responses attributable to distinct identities
		is therefore at most $m \cdot \tau_h$.
	\end{proof}
	
	\begin{theorem}[Linear Sybil cost]
		\label{thm:linear-cost}
		Under~(H1)--(H4), sustaining $s$ active identities in any single
		window requires $m = \Omega(s)$ independent participation channels.
		Equivalently, $C(s, 1) = \Omega(s)$.
	\end{theorem}
	
	\begin{proof}
		From Lemma~\ref{lem:identity-bound}, $m\tau_h \ge s$, giving
		$m \ge s/\tau_h = \Omega(s)$ since $\tau_h = O(1)$. Each channel-window
		incurs a non-zero participation cost, so $C(s,1) = \Omega(s)$.
	\end{proof}
	
	\begin{corollary}[No sublinear amortization]
		No adversary can sustain $s$ identities with $o(s)$ channels per window.
	\end{corollary}
	
	\begin{proof}
		Amortizing one channel across multiple identities in the same window
		violates~(H1). Cross-window reuse violates~(H2)--(H3). Both are
		excluded by construction.
	\end{proof}
	
	
	%% =============================================================
	\section{Multi-Window and Long-Term Analysis}
	\label{sec:multiwindow}
	%% =============================================================
	
	Let $s_d$ and $m_d$ denote the number of active identities and channels,
	respectively, in window $d$.
	
	\begin{lemma}[Window persistence bound]
		\label{lem:persistence}
		$s_d \le m_d \cdot \tau_h$ for every window $d$.
	\end{lemma}
	
	\begin{proof}
		Lemma~\ref{lem:identity-bound} applied to window $d$. Since~(H2)--(H3)
		prevent cross-window reuse, effort in window $d$ cannot carry to $d+1$.
	\end{proof}
	
	\begin{theorem}[Joint linear cost]
		\label{thm:joint-linear}
		Under~(H1)--(H4), $C(s,T) = \Omega(sT)$.
	\end{theorem}
	
	\begin{proof}
		By Lemma~\ref{lem:persistence}, each of $T$ windows independently
		requires $\Omega(s)$ channels. Channels cannot be reused across
		windows~(H2)--(H3), so the total cost is
		$\sum_{d=1}^{T}\Omega(s) = \Omega(sT)$.
	\end{proof}
	
	\begin{theorem}[Steady-state Sybil capacity]
		\label{thm:steady-state}
		Let $\bar{m} = \limsup_{T\to\infty}\frac{1}{T}\sum_{d=1}^{T}m_d$.
		Then $\;\limsup_{T\to\infty}\frac{1}{T}\sum_{d=1}^{T}s_d
		\;\le\; \bar{m}\cdot\tau_h$.
	\end{theorem}
	
	\begin{proof}
		Sum $s_d \le m_d\tau_h$ over $T$ windows, divide by $T$, take
		$\limsup$.
	\end{proof}
	
	\noindent\textbf{Burst attacks.}
	Concentrating effort in short bursts does not help: the total
	identity-window budget $B_T = \sum_{d=1}^T s_d$ satisfies
	$B_T \le \tau_h\sum_{d=1}^T m_d$, so burst attacks incur linear
	cost in total channel-windows regardless of temporal concentration.
	
	
	%% =============================================================
	\section{Challenge Families}
	\label{sec:challenges}
	%% =============================================================
	
	\subsection{Abstract Admissibility Requirements}
	
	A challenge family $\mathcal{C}$ is \emph{admissible} for \HCO\ if
	every instance satisfies~(H1)--(H4). Concretely:
	\begin{itemize}
		\item instances are unpredictable before issuance~(H2);
		\item correct responses are deterministically verifiable;
		\item responses are cryptographically bound to $(v,d,j)$~(H1)
		via the construction of Section~\ref{sec:construction};
		\item the interaction model imposes an irreducible per-response
		latency that limits channel throughput to $\tau_h$ per window~(H4).
	\end{itemize}
	The construction of Section~\ref{sec:construction} satisfies~(H1)--(H3)
	for any task payload; admissibility under~(H4) depends on the interaction
	latency of the specific task family.
	
	\subsection{Admissible Families}
	
	\noindent\textbf{Perceptual visual matching.}
	A distorted or noise-corrupted image is shown with candidate matches;
	the participant selects the correct correspondence within
	$\Delta_{\mathrm{resp}}$. Human perceptual generalization under noise
	remains robust~\cite{geirhos2018generalisation}; current VLMs incur
	multi-second inference latencies that enforce a natural throughput ceiling
	per automated channel. Instances are generated by random visual
	transformations; verification is deterministic.
	
	\noindent\textbf{Interactive reasoning tasks.}
	Short dependent reasoning sequences require maintaining intermediate
	state and submitting a result before a visible countdown expires.
	Multi-step inference latency limits automated throughput. Correctness
	is verified deterministically from dynamically generated parameters.
	
	\noindent\textbf{Biometric-light response tasks.}
	A fresh, unpredictable prompt requires a synchronized real-time signal
	(spoken phrase, motion gesture). Automated synthesis under real-time
	constraints is currently unreliable and produces detectable latency
	artifacts. Verification uses transient signal features without
	long-term biometric enrollment.
	
	\noindent\textbf{Attention-based interaction tasks.}
	A dynamic scene requires continuous tracking or selection over several
	seconds, occupying the full $\Delta_{\mathrm{resp}}$ interval. This
	structurally limits any channel to one response per window.
	
	\subsection{Deployment Example: Relaxed Window Sizes}
	\label{sec:deployment-example}
	
	The formal guarantees of~\HCO\ hold across the full parameter space of
	$\Delta_{\mathrm{resp}}$ and $\tau_h$, from sub-second strict deadlines
	to relaxed daily interaction windows. A concrete example illustrates
	the range.
	
	\noindent\textbf{Daily participation with device-bound login.}
	Consider a system in which each identity must perform one login-style
	interaction per day: $\Delta_h = 24\,\text{h}$, $\Delta_{\mathrm{resp}}
	= 24\,\text{h}$, and $\tau_h = 1$. A single human has 24 hours to
	respond---there is no cognitive pressure. However, if the login is
	\emph{device-bound} (bound to a specific smartphone, secure element, or
	attested browser session), then sustaining $s$ identities requires $s$
	independent attested devices or sessions. An adversary cannot amortize
	one device across $s$ identities within the same window: each identity
	requires its own channel-window.
	
	Under this configuration, Theorem~\ref{thm:joint-linear} gives
	$C(s,T) = \Omega(sT)$ channel-windows even with $\tau_h = 1$. Sybil
	resistance does not depend on the challenge being cognitively demanding
	or time-pressured; it depends entirely on channel non-transferability
	(H1) and window-locality~(H4) enforced by the device binding.
	
	This example also clarifies the design space: tighter
	$\Delta_{\mathrm{resp}}$ values (seconds to minutes) provide stronger
	near-real-time resistance but impose higher usability cost. Relaxed
	windows (hours to daily) reduce user burden while preserving the
	linear cost structure, provided channel binding is enforced
	engineering-side. The appropriate point in this tradeoff is a
	deployment-level decision orthogonal to the formal guarantees.
	
	\subsection{Rotation and Composability}
	
	\HCO\ does not depend on any single family. Systems may rotate across
	admissible families between windows or combine them within a window.
	Since security depends only on~(H1)--(H4), rotation preserves all
	guarantees of Theorems~\ref{thm:joint-linear}
	and~\ref{thm:steady-state}---this is essential for long-lived
	deployments where the automation landscape evolves.
	
	
	%% =============================================================
	\section{Empirical Evaluation}
	\label{sec:eval}
	%% =============================================================
	
	We conduct a preliminary empirical study to validate~(H4): that
	representative challenge families impose a measurable throughput bound
	on current automated systems under strict response deadlines, and that
	honest human participants complete challenges at high rates within
	$\Delta_{\mathrm{resp}}$. We emphasize that these results are preliminary;
	a large-scale study is planned as future work
	(see Section~\ref{sec:limitations}).
	
	\subsection{Methodology}
	
	\noindent\textbf{Human participants.}
	We recruited 30 participants (ages 18--45) through university mailing
	lists. Each participant completed 20 challenges per family (80 total)
	via a browser-based interface with a visible countdown timer.
	
	\noindent\textbf{Automated solvers.}
	We evaluated GPT-4o~\cite{openai2023gpt4}, Gemini~1.5~Pro~\cite{geminiteam2023},
	and Claude~3.5~Sonnet~\cite{anthropic2024claude} via their standard
	vision-language inference APIs, plus a specialized CAPTCHA-solving
	service (2captcha) for the perceptual family. We ran 100 trials per
	(solver, family) pair. A trial was counted as successful only if a
	correct response was \emph{received} within the simulated
	$\Delta_{\mathrm{resp}}$ window; inference latency exceeding
	$\Delta_{\mathrm{resp}}$ counted as failure regardless of correctness.
	
	\subsection{Results}
	
	\begin{table*}[t]
		\centering
		\caption{Preliminary results under strict $\Delta_{\mathrm{resp}}$
			constraints. Automated success includes both incorrect
			solutions and deadline violations.}
		\label{tab:evaluation-results}
		\resizebox{\linewidth}{!}{%
			\begin{tabular}{lcccc}
				\toprule
				Challenge Family
				& Human Success (\%)
				& Human Mean Time (s)
				& Automated Success (\%)
				& Auto.\ Mean Time (s) \\
				\midrule
				Perceptual visual matching  & 92  & 6.2  & 12          & 18.4 \\
				Interactive reasoning       & 85  & 11.8 & 18          & 22.1 \\
				Biometric-light response    & 100 & 8.1  & $\approx 0$ & --   \\
				Attention-based interaction & 95  & 14.5 & $\approx 0$ & --   \\
				\bottomrule
			\end{tabular}%
		}
	\end{table*}
	
	Human participants achieved high success rates across all families,
	with mean completion times well below typical $\Delta_{\mathrm{resp}}$
	values. Automated solvers were primarily limited by inference latency:
	mean response times of 18.4\,s and 22.1\,s for perceptual and
	reasoning families substantially exceed the 8--15\,s deadlines tested,
	confirming that current VLMs impose a natural throughput ceiling
	consistent with~(H4).
	
	
	%% =============================================================
	\section{Security Analysis}
	\label{sec:security}
	%% =============================================================
	
	We analyze \HCO\ against the principal adversarial strategies. All
	arguments are grounded in~(H1)--(H4) and the construction of
	Section~\ref{sec:construction}.
	
	\noindent\textbf{Automated solving.}
	An adversary deploying AI solvers to fill participation channels does
	not circumvent the linear cost bound. By~(H4), each automated channel
	is throughput-bounded at $\tau_h$ responses per window under
	$\Delta_{\mathrm{resp}}$. Scaling to $s$ identities requires $\Omega(s)$
	independent channels. Automation shifts the cost from human labor to
	API credits or hardware, but preserves linear scaling in $s$.
	
	\noindent\textbf{Human outsourcing and labor markets.}
	An adversary may hire human workers or use CAPTCHA-solving
	services~\cite{motoyama2010recaptchas} to solve challenges.
	This attack is captured in the adversary model: the adversary operates
	$m$ human channels. Sustaining $s$ identities requires $m = \Omega(s)$
	workers (Theorem~\ref{thm:linear-cost}). Outsourcing converts the cost
	from direct labor to wages, but the linear scaling is preserved.
	
	\noindent\textbf{Relay and real-time forwarding attacks.}
	Relaying challenges to external solvers adds network round-trip time,
	which tightens rather than relaxes~(H4). By~(H1) and~(H2)--(H3),
	relayed responses are identity-bound and non-reusable; relaying does
	not reduce the total independent channel-windows required.
	
	\noindent\textbf{Replay and solution reuse attacks.}
	By~(H1), response $\rho_{v,d,j}$ is bound to $(v,d,j)$ via the
	signature $\sigma_v$ and the binding tag. Submitting the same response
	for $(v',d',j') \neq (v,d,j)$ fails the binding-tag check in
	Algorithm~\ref{alg:verify} (lines~3--5). By~(H2), the nonce $\eta_d$
	ensures responses prepared before window $d$ are invalid.
	
	\noindent\textbf{Precomputation and batch-solving.}
	By~(H2), $\chi_{v,d,j}$ cannot be computed before $\eta_d$ is revealed.
	Batch-solving within a window is bounded by~(H4): processing $k > \tau_h$
	challenges requires more than $\tau_h$ round-trip latency budgets,
	which is infeasible under $\Delta_{\mathrm{resp}}$.
	
	\noindent\textbf{Synthetic media and deepfake attacks.}
	For biometric-light families, an adversary may attempt real-time synthesis
	of a biometric response. Mitigations include: (i)~challenge freshness~(H2)
	prevents off-line synthesis; (ii)~liveness detection flags synthesized
	signals; (iii)~challenge rotation to families resistant to synthesis
	(e.g., attention-based interaction) eliminates reliance on vulnerable
	modalities.
	
	\noindent\textbf{Side-channel and man-in-the-browser attacks.}
	A malicious browser extension may observe and forward challenges to an
	external solver. This constitutes a relay attack and is bounded as above.
	Deployment-level countermeasures---content-security-policy headers,
	subresource integrity, and challenge obfuscation---reduce exfiltration
	feasibility but are orthogonal to the core security model.
	
	\noindent\textbf{Privacy and data minimization.}
	Verification requires only $\rho = (\sigma_v, r_{v,d,j})$ and
	$\chi_{v,d,j}$; no raw biometric data or behavioral record is stored.
	$\mathsf{pk}_v$ may be pseudonymous. Privacy-preserving challenge
	delivery ensuring unlinkability across windows while preserving identity
	binding is an important direction for future work.
	
	
	%% =============================================================
	\section{Limitations}
	\label{sec:limitations}
	%% =============================================================
	
	\noindent\textbf{Dependence on Property~(H4).}
	The linear cost guarantee rests critically on~(H4). For human channels,
	this bound is grounded in cognitive and temporal constraints. For
	automated channels, it is grounded in inference and network latency.
	Advances in specialized hardware or ultra-low-latency inference pipelines
	could reduce the effective latency floor, increasing $\tau_h$ and
	weakening the guarantee. Long-lived deployments must periodically
	reassess~(H4) and adjust $\Delta_{\mathrm{resp}}$ accordingly, analogous
	to difficulty adjustment in PoW~\cite{nakamoto2008bitcoin}.
	
	\noindent\textbf{Engineering bypass and the nature of~(H4).}
	When~(H4) is enforced through channel serialization rather than
	cognitive constraints, adversaries may attempt to scale channel capacity
	through infrastructure: VM farms, cloud phone emulators, browser-bot
	fleets, or distributed worker pools. Such infrastructure does not
	invalidate~(H4); it instead converts the guarantee from an
	\emph{impossibility} into a \emph{cost multiplier}. Each additional
	sustained identity still requires provisioning and operating an
	independent channel, so the adversarial cost remains $\Omega(sT)$---but
	$\tau_h$ and the per-channel cost are now determined by engineering
	practice rather than cognition.
	
	This is analogous to the role of honest-majority in Byzantine consensus:
	the guarantee holds given the assumption, and the assumption has a
	concrete cost to violate. The honest disclosure of this structure---that
	\HCO\ bounds adversarial \emph{cost} rather than adversarial
	\emph{capability}---makes the primitive more trustworthy, not less.
	Deployments relying on channel serialization should therefore size
	$\Delta_{\mathrm{resp}}$ and per-channel costs such that operating a
	large farm is economically prohibitive for the target threat model.
	
	\noindent\textbf{Scope of the preliminary evaluation.}
	The study uses 30 participants recruited via university mailing lists,
	not a controlled recruitment platform~\cite{prolific2014,peer2022data}.
	It does not cover task-specific fine-tuned models, custom hardware
	solvers, or organized human--AI hybrid pipelines. A large-scale study
	(N~$\geq$~100 via Prolific, comprehensive adversarial solver suite,
	statistical significance testing) is planned as future work.
	
	\noindent\textbf{Accessibility and inclusivity.}
	Challenge families requiring visual, auditory, or motor interaction may
	exclude participants with relevant impairments. Deployments must provide
	alternative modalities and accessibility-aware selection policies.
	
	\noindent\textbf{Parameter sensitivity.}
	$\Delta_{\mathrm{resp}}$ and $k(d)$ jointly determine the security-usability
	tradeoff. The formal guarantees hold for any $\Delta_{\mathrm{resp}} > 0$,
	but concrete parameterization requires empirical measurement of $\tau_h$
	and honest participant success rates in the target deployment environment.
	
	\noindent\textbf{Theoretical scope.}
	The cost analysis assumes~(H1)--(H4) hold throughout execution. Partial
	failures of~(H4) yield intermediate scaling regimes; their characterization
	is developed in the companion paper~\cite{maleki2026geometry}.
	
	
	%% =============================================================
	\section{Conclusion}
	\label{sec:conclusion}
	%% =============================================================
	
	We introduced the Human Challenge Oracle~(\HCO), a cryptographic
	primitive for continuous, identity-bound, rate-limited participation
	verification. Under Properties~(H1)--(H4), \HCO\ enforces
	$C(s,T) = \Omega(sT)$: sustaining $s$ identities over $T$ windows
	requires $\Omega(sT)$ independent channel-windows, regardless of
	whether channels are operated by humans, automated agents, or any
	combination thereof. No assumption about human cognitive superiority
	is required.
	
	We presented a hash-based construction satisfying~(H1)--(H3) under
	standard cryptographic assumptions, with Algorithm~\ref{alg:verify}
	as the formal verification procedure. We identified four admissible
	challenge families, described browser-based instantiations, and
	conducted a preliminary empirical study confirming that current VLMs
	impose a natural throughput ceiling under strict response deadlines,
	validating the empirical basis of~(H4).
	
	\HCO\ is designed as a composable primitive. Its formal use as a
	consensus resource---including backbone-style safety, liveness, and
	fairness---is developed in~\cite{maleki2026pocmt}. The axiomatic
	cost-theoretic foundation establishing that linear Sybil cost requires
	a non-parallelizable resource is given in~\cite{maleki2026geometry}.
	
	The most immediate priorities for future work are a large-scale empirical
	study (N~$\geq$~100, controlled recruitment, comprehensive adversarial
	solvers), privacy-preserving challenge delivery with cross-window
	unlinkability, formal composability analysis in the UC framework, and
	adaptive challenge rotation mechanisms that preserve~(H1)--(H4) under
	evolving automation capabilities.
	
	
	\bibliographystyle{unsrtnat}
	\bibliography{references}
	
	
\end{document}